\section{Euler references, density and exceptional centers}
\label{sec:euler-reference}

A finite Euler reference has two advantages over a fixed arithmetic
prefix: its logarithm calibrates exactly to a prime subseries, and
its modulus never vanishes. The reflected reference does have real
zeros; these will be retained inside integrable logarithms. The
development combines E05's fixed-reference moment analysis and E06's
joint-phase estimate \cite{GNPRSEuler,GNPRSAdaptive}.

\subsection{Exact calibration and a regulator-uniform phase lemma}
For \(P\ge2\), fixed across the spatial integral, define
\begin{equation}\label{eq:euler-reference}
 E_P(x)=\prod_{p\le P}(1-p^{-\sigma_c}
                   \ee^{\ii\nu_\lambda(x)\log p})^{-1},\qquad
 \sigma_c=\tfrac12+c/4,\qquad
 S_P^{\rm ref}=E_P+\ee^{\ii\phi_\lambda}\overline E_P.
\end{equation}
Use the canonical logarithm
\(\ell_P=\sum_{p\le P}\sum_{k\ge1}
 p^{-k\sigma_c}\ee^{\ii k\nu_\lambda\log p}/k\).
The finite product and its inverse are bounded on the entire real
line for fixed \(P\). Absolute convergence and the Fourier transform
of \(K\) give
\begin{equation}\label{eq:euler-calibration}
 \int_\R K(x-X)\log|E_P(x)|\dd x
 =P_{c,\lambda}^{(\le P)}(X)
 :=\sum_{p\le P}\sum_{k\ge1}
 \frac{p^{-k(5/2+c/4)}\cos(k\nu_\lambda(X)\log p)}{k(1+p^{-4k})}.
\end{equation}
All powers of the retained primes are included. The omitted prime
subseries costs at most \((2/3)P^{-3/2}\), since every omitted
prime power exceeds \(P\).

Put \(\kappa_E=1-2^{-5/8}>0\). Differentiating the absolutely
convergent logarithm, using \(\nu_\lambda'=1/2\), yields
\begin{align}\label{eq:euler-phase-derivatives}
 2|\Im\ell_P'|&\le D_1(P):=\frac8{3\kappa_E}P^{3/8}\log P,\nonumber\\
 2|\Im\ell_P''|&\le D_2(P):=\frac4{3\kappa_E^2}P^{3/8}(\log P)^2.
\end{align}
For example the first sum is bounded by
\(\sum_{p\le P}p^{-5/8}\log p/(1-p^{-5/8})\); the second has
coefficient \(1/2\) and denominator squared. In both, use
\(\sum_{n\le P}n^{-5/8}\le(8/3)P^{3/8}\).

\begin{lemma}[The complete reference phase]\label{lem:joint-phase}
Let \(\psi=\phi_\lambda-2\Im\ell_P\), so that
\(a:=S_P^{\rm ref}/E_P=1+\ee^{\ii\psi}\). Assume the eventual
window gates and \(D_1(P)\le L_X/8\), \(R_X/X\le L_X/8\).
Then \(\psi'\ge L_X/4\) on \(W_X\), and, writing
\(D=\|\psi''\|_{\infty,W_X}\),
\begin{align}\label{eq:joint-phase}
 D&\le D_2(P)+X^{-1}+\pi\lambda/(2X^2),\nonumber\\
 \left|\int_{W_X}\tfrac12\log(|a|^2+\delta^2)\,d\omega_X
                   -q_X\log b_\delta\right|
 &\le\frac2{L_X}+\frac{16D}{L_X^2},
 \qquad b_\delta=\frac{\delta+\sqrt{\delta^2+4}}2.
\end{align}
The bound holds uniformly for every \(\delta>0\) fixed within
the integral, and also at \(\delta=0\), with \(\log b_0=0\).
\end{lemma}
\begin{proof}
At fixed time \(\phi_\lambda'=L_x/2+\pi\lambda/(8x)\) and
\(\phi_\lambda''=1/(2x)-\pi\lambda/(8x^2)\). The derivative
bounds and \(L_x\ge L_X-2R_X/X\) prove the assertions about \(\psi\).
The unit-arm instance of \eqref{eq:regularized-phase-factor} is
\[
 \tfrac12\log(|a|^2+\delta^2)=\log b_\delta+
  \Re\sum_{k\ge1}\frac{(-1)^{k+1}\rho^k}{k}\ee^{\ii k\psi},
 \qquad \rho=b_\delta^{-2}<1.
\]
Its coefficient modulus is constant in \(x\). Integration by parts
in the entire phase, with \(m=\min\psi'\), gives
\[
 \left|\int_{W_X}\ee^{\ii k\psi}\,d\omega_X\right|
 \le\frac{2\|K\|_\infty+\|K'\|_1}{km}
          +\frac{q_XD}{km^2}.
\]
Here \(\|K\|_\infty=1/16\), \(\|K'\|_1=1/8\) and
\(q_X\le1/2\). Summing with \(\rho^k/k\) and
\(\sum\rho^k/k^2\le2\) proves \eqref{eq:joint-phase}.
For \(0<\delta\le1\), the regularized log decreases to
\(\log|a|\) and is bounded above by \(\tfrac12\log5\).
Monotone convergence applied to this constant minus the log, and
the uniform two-sided integral bound, give a finite limit with the
same estimate. Thus the real reference zeros are included.
\end{proof}

Define the exact relative residual
\begin{equation}\label{eq:exact-relative-residual}
 R_P(X,c)=\int_{W_X}\frac{|V_\lambda-S_P^{\rm ref}|^2}{|E_P|^2}
                   \,d\omega_X.
\end{equation}
Only the unreflected nonvanishing product is divided out. Also
\(|\ell_P|\le(8/(3\kappa_E))P^{3/8}\), so its omitted exterior
costs at most
\(B_P^{\rm ext}X^{-4}\), where
\(B_P^{\rm ext}=16P^{3/8}/(3\pi\kappa_E)\).
Applying Lemma~\ref{lem:upper-log-transfer} to \(V/E_P,a\), then
Jensen only to the residual payment, proves the complete bound
\begin{align}\label{eq:euler-master}
 J_\lambda(X)\le{}&P_{c,\lambda}^{(\le P)}-P_{c,\lambda}
       +q_X\log b_\delta
       +\frac{q_X}{2}\log\left(1+\frac{R_P(X,c)}{q_X\delta^2}\right)
       \nonumber\\
 &+T_X+\frac{C_{\rm out}+2}{L_X}
       +\frac{16D}{L_X^2}+B_P^{\rm ext}X^{-4}.
\end{align}
No field floor and no omitted true logarithmic exterior enter this
inequality. For fixed \(P\), its phase error is \(O_P(1/L_X)\)
uniformly over regulators, strengthening the older fixed-regulator
form of E05 by E06's phase lemma.

\subsection{The fixed-reference residual mean and density gain}
\begin{theorem}[Fixed-prime residual law]\label{thm:residual-mean}
For each fixed finite \(P\),
\begin{equation}\label{eq:residual-mean}
 \lim_{T\to\infty}\frac1T\int_T^{2T}R_P(X,1)\dd X
 =\mu_P:=\zeta(3/2)\prod_{p\le P}(1-p^{-3/2})-1,
 \qquad \mu_P\longrightarrow0\quad(P\to\infty).
\end{equation}
The height limit is taken before the prime-cutoff limit.
\end{theorem}
\begin{proof}
At the auxiliary time \(1/L_u\), let \(D_P=\prod_{p\le P}p\) and
\[
 Q(u)=E_P(u)^{-1}=\sum_{d\mid D_P}\mu(d)d^{-3/4}
                  \ee^{\ii\nu(u)\log d},\qquad B(u)=A(u)Q(u)-1,
\]
where \(\mu(d)\) is the M\"obius function. The coefficient at 1
vanishes. For \(k>1\), the coefficient of
\(\ee^{\ii\nu(u)\log k}\) is
\[
 b_k(u)=\sum_{\substack{d\mid D_P,\ d\mid k\\ k/d\le N_u}}
                     \mu(d)d^{-3/4}w_{k/d}(u).
\]
Its supremum plus total variation on \([T,2T]\) is
\(O_P(k^{-5/8})\), including each cutoff activation. For fixed
\(k>1\) it tends uniformly to \(k^{-3/4}\mathbf1_{(k,D_P)=1}\).
The diagonal thus tends to
\(\sum_{k>1,(k,D_P)=1}k^{-3/2}=\mu_P\), by the summable square
majorant \(O_P(k^{-5/4})\). The same-arm argument in
Theorem~\ref{thm:coherent-moment}, now through \(D_PN_{2T}\),
pays all off-diagonal terms by \(O_P(T^{-5/8}(1+\log T))\).

The reflected relative square equals
\[
 2|B|^2+2\Re\{\ee^{-\ii\phi}B^2\overline Q/Q\}.
\]
The finite-prime function \(\overline Q/Q=\overline Q E_P\) has
an absolutely uniformly convergent expansion in rational frequencies
\(\ee^{\ii\nu\log r}\). Truncate it with uniform error
\(\epsilon\). For each of the finitely retained \(r\), the phase
\(-\phi+\nu\log(klr)\) has second derivative of fixed sign and
magnitude \(\asymp1/T\), since \(k,l\ll_P\sqrt T\).
The previous second-derivative argument costs \(O_{P,r}(T^{-1/8})\)
after summation. The removed frequencies cost at most
\(2\epsilon\) times the bounded mean of \(|B|^2\).
Take \(T\to\infty\), then \(\epsilon\downarrow0\); the cross
mean is zero.

Dividing the complete field error by \(E_P\) costs
\(O_P(u^{-1/4})\), hence \(O_P(u^{-1/16})\) in the square.
The auxiliary mean is therefore \(2\mu_P\).
The fixed-clock comparison has the same
\(O_P(R_XX^{-9/16}+X^{-1/32})\) bound as before, because the fixed
product and its inverse have bounded modulus and time derivative;
the reflected phase derivative is already included in the arm bound.
The same Fubini edges and kernel mass \(1/2\) prove
\eqref{eq:residual-mean}. Finally \(\mu_P\to0\) by domination in
the displayed coprime sum.
\end{proof}

\begin{corollary}[Density of centers with a large Green value]
\label{cor:green-density}
For every fixed \(h>0\),
\begin{equation}\label{eq:green-density}
 \lim_{T\to\infty}\frac1T
   |\{X\in[T,2T]:J_{1/L_X}(X)\ge h\}|=0.
\end{equation}
In particular this holds at the original defect margin \(h=m_*\).
There is no rate or assertion about a fixed time, all \(c\), or
the density of zeros.
\end{corollary}
\begin{proof}
Choose a fixed \(\delta>0\) so \(q_X\log b_\delta\le h/5\),
and then arbitrarily large fixed \(P\) so the prime tail is at most
\(h/5\). Set \(r=\delta^2h/5\). If \(R_P\le r\), its
logarithmic payment is at most \(r/(2\delta^2)=h/10\).
The remaining errors in \eqref{eq:euler-master} are eventually at
most \(h/5\) for each fixed \(P\). Hence the bad set is contained
in \(\{R_P>r\}\). Markov's inequality and
Theorem~\ref{thm:residual-mean} bound its upper dyadic density by
\(\mu_P/r\). Let fixed \(P\) tend to infinity after the height
limit. This proves the statement for arbitrary \(h\), explicitly
extracting the arbitrary-threshold consequence of E05's proof at
\(m_*\).
\end{proof}

\subsection{An averaged consequence for the actual Green value}
The density conclusion can be strengthened in mean by the coherent
moment bound. The following consequence is derived in this synthesis
\cite{GNPRSMeanConsequences}; it adds no exceptional-center arithmetic
estimate to E05.

\begin{theorem}[Power means and subcritical exponential means]
\label{thm:green-power-means}
Write \(J(X)=J_{1/L_X}(X)\). For all sufficiently large \(T\),
\begin{equation}\label{eq:green-exponential-bound}
 \frac1T\int_T^{2T}\ee^{4J(X)}\dd X\le39.
\end{equation}
For every fixed \(p>0\) and every fixed \(0<\alpha<4\),
\begin{equation}\label{eq:green-power-means}
 \frac1T\int_T^{2T}J(X)^p\dd X\longrightarrow0,
 \qquad
 \frac1T\int_T^{2T}\ee^{\alpha J(X)}\dd X\longrightarrow1.
\end{equation}
The constant 39 is an explicit convenient bound; its finite onset is
unevaluated. Neither convergence rate nor convergence at \(\alpha=4\)
is asserted.
\end{theorem}
\begin{proof}
Eventually \(q_X\ge1/4\) and
\(T_X+C_{\rm out}/L_X\le1/100\). Using
\(|P_{1,1/L_X}|<23/100\) in \eqref{eq:quadratic-allowance} gives
\[
 \ee^{4J(X)}\le\ee^{24/25}
       \left(\frac{M(X,1)}{q_X}\right)^{2q_X}
       \le\ee\,(1+4M(X,1)).
\]
The last inequality uses \(t^r\le1+t\) for \(t\ge0\),
\(0<r\le1\). Also
\(\zeta(3/2)\le1+2^{-3/2}+\int_2^\infty u^{-3/2}\dd u<3\).
Theorem~\ref{thm:coherent-moment} thus makes the averaged right side
at most \(13\ee<39\) eventually.

Let \(d_T(t)=T^{-1}|\{X\in[T,2T]:J(X)>t\}|\).
Positivity of the Green representation and Markov's inequality give
\(d_T(t)\le\min(1,39\ee^{-4t})\). For each fixed \(t>0\),
Corollary~\ref{cor:green-density} gives \(d_T(t)\to0\).
Tonelli's theorem yields
\[
 \frac1T\int_T^{2T}J^p\dd X
     =p\int_0^\infty t^{p-1}d_T(t)\dd t,
 \qquad
 \frac1T\int_T^{2T}(\ee^{\alpha J}-1)\dd X
     =\alpha\int_0^\infty\ee^{\alpha t}d_T(t)\dd t.
\]
The displayed tail bound is an integrable majorant in the first
identity for every \(p>0\), and in the second for \(\alpha<4\).
Dominated convergence proves \eqref{eq:green-power-means}.
\end{proof}

In particular the actual Green potential tends to zero in every
finite \(L^p\) mean on dyadic intervals along this moving-time slice.
Narrow exceptional sets can survive every such mean statement.
There is no assertion of pointwise convergence or a supremum bound.
For an abstract limit check on a unit probability interval, take
\(Y_n=(\log n)\mathbf1_{[0,1/n]}/4\). Then every power mean
tends to zero and every exponential mean of exponent less than four
tends to one, but \(\mathbb E\ee^{4Y_n}=2-1/n\) and
\(\|Y_n\|_\infty\to\infty\). This is a measure-theoretic
example, not a theta field or a zero construction; it identifies the
logical limit of the mean statements alone.

\subsection{Exceptional centers for every fixed reference}
The mean law is compatible with unbounded exceptional residuals.
The reason is arithmetic coherence, which can be detected by fixed
rational-frequency tests without any quantitative equidistribution
rate.

\begin{lemma}[Twisted complete moments]\label{lem:twisted-moments}
For fixed coprime positive integers \(a,b\), let \(s=\log(a/b)\).
Then
\begin{equation}\label{eq:twisted-moment}
 \lim_{T\to\infty}\frac1T\int_T^{2T}\ee^{\ii Xs/2}M(X,1)\dd X
     =2\zeta(3/2)(ab)^{-3/4}\widehat K(s/2).
\end{equation}
\end{lemma}
\begin{proof}
In the auxiliary same-arm square the phase is
\(\nu(u)\log(n/m)+us/2\). Resonance requires \(an=bm\),
so \(n=bj,m=aj\), and its coefficient sum tends to
\((ab)^{-3/4}\zeta(3/2)\) for each of the two arm squares.
For nonresonant pairs
\(|\log(an/(bm))|\gg_{a,b}T^{-1/2}\).
The \(O(1/(T\log T))\) derivative perturbation preserves that
lower bound; the monotone first-derivative estimate costs
\(O_{a,b}(\sqrt T)\) per partial integral. The summed normalized
error is \(O_{a,b}(T^{-1/8})\). The linear twist leaves the
cross-arm second derivative unchanged. The correction and fixed-clock
payments are unchanged in absolute value. In the Fubini step the
twist contributes \(\ee^{-\ii ts/2}\) to the kernel at \(t=x-X\),
giving \(\widehat K(s/2)\) in place of its mass. The same edges
and tails finish the proof.
\end{proof}

\begin{theorem}[Fixed-reference exceptional residuals]
\label{thm:fixed-residual-unbounded}
The exact moment \(M(X,1)\) is unbounded at arbitrarily high centers.
For every fixed \(P\), so is \(R_P(X,1)\).
\end{theorem}
\begin{proof}
Choose a second fixed prime cutoff \(U\), independent of \(P\).
On its finite torus, with probability Haar measure \(dz\), put
\[
 E_U(z)=\prod_{p\le U}(1-p^{-3/4}z_p)^{-1},\quad
 C_U=\zeta(3/2)\prod_{p\le U}(1-p^{-3/2}),\quad
 \tau_p(t)=\ee^{\ii t\log p/2},
\]
\[
 F_U(z)=2C_U\int_\R K(t)|E_U(z\tau(t))|^2\dd t.
\]
This is continuous. For a character corresponding to
\(a/b=\prod p^{\ell_p}\), expansion of the finite-prime product
shows that \(\int z^\ell F_U(z)\dd z\) is the right side of
\eqref{eq:twisted-moment}. Along \(z_p(X)=\ee^{\ii X\log p/2}\),
every nonconstant character has zero limiting height mean, by unique
factorization. Thus each fixed trigonometric polynomial has the
corresponding weighted height mean against \(M\).

Use the nonnegative product Fej\'er polynomials
\[
 W_N(z)=\prod_{p\le U}\frac1N\left|\sum_{j=0}^{N-1}z_p^j\right|^2.
\]
Their Haar mean is one and they concentrate at \(z=1\): outside
\(|\arg z_p|<\eta\), the relevant factor is at most
\(1/[N\sin^2(\eta/2)]\). A finite union bound and continuity give
\(\int W_NF_U\dd z\to F_U(1)\).
If \(M\le H\) eventually, take the height limit first for fixed
\(U,N\) in the nonnegative inequality \(W_NM\le HW_N\), then
let \(N\to\infty\). This would imply \(F_U(1)\le H\) for every
fixed \(U\).

But for \(|t|\le1/\log U\) the prime cosines are at least \(7/8\),
and
\[
 |E_U(\tau(t))|^2\ge
        \exp\left(\frac34\sum_{p\le U}p^{-3/4}\right).
\]
The harmonic sum through \(\lfloor U\rfloor\) is at most
\(\prod_{p\le U}(1-1/p)^{-1}\); taking logarithms and paying the
summable higher powers gives
\(\sum_{p\le U}1/p\ge\log\log U-O(1)\).
Since \(p^{1/4}\ge4\) above a fixed cutoff,
\(\sum_{p\le U}p^{-3/4}\ge4\log\log U-O(1)\).
As \(C_U\ge1\) and \(K(t)\ge K(1)\) on this short interval,
\(F_U(1)\gg(\log U)^2\), a contradiction.

Finally \(B_P:=\sup_x|E_P(x)|<\infty\) for fixed \(P\). The
weighted triangle inequality gives
\begin{equation}\label{eq:residual-unbounded-transfer}
 R_P(X,1)\ge B_P^{-2}
        \bigl(\sqrt{M(X,1)}-2B_P\sqrt{q_X}\bigr)_+^2,
\end{equation}
which proves the residual assertion.
\end{proof}

This also obstructs the specified separate residual/Jensen payment
even with a regulator chosen at each center. Before any phase estimate,
the regularized reference log is at least \(q_X\log\delta\); hence
its sum with the residual payment is at least
\[
 \frac{q_X}{2}\log(\delta^2+R_P/q_X)
      \ge\frac{q_X}{2}\log(R_P/q_X)
\]
when \(R_P>0\). Restoring the bounded fixed-reference prime
difference cannot prevent divergence at those centers. This is a
statement about that allowance, not about the actual Green value.

\subsection{A growing reference with paid comparison costs}
Now choose \(P=P(X)=\lfloor L_X^2\rfloor\), retaining this one
value throughout each window. Equations~\eqref{eq:euler-phase-derivatives}
give \(D_1(P),D_2(P)=O(L_X^{3/4}(\log L_X)^2)=o(L_X)\).
Thus the phase gates hold and \(D\le L_X\) eventually, uniformly
in \(c\). The joint phase costs at most \(18/L_X\), for every
positive regulator. More generally the same conclusion holds for
\(2\le P\le L_X^\gamma\), with any fixed \(0<\gamma<8/3\).
At the selected cutoff the prime tail is \(O(L_X^{-3})\) and the
reference-log exterior is \(O(L_X^{3/4}X^{-4})\).

The full correction and provider error remain small after normalization:
\begin{equation}\label{eq:adaptive-provider-error}
 \sup_{W_X}\left|\frac{V_\lambda-S_\lambda}{E_P}\right|
 \le CX^{-1/4}\exp\{(8/3)L_X^{3/4}\}=X^{-1/4+o(1)}.
\end{equation}
Indeed \(|E_P|^{-1}\le\prod_{p\le P}(1+p^{-\sigma_c})
\le\exp((8/3)P^{3/8})\). The complete adaptive inequality is
\begin{align}\label{eq:adaptive-allowance}
 J_\lambda(X)\le{}&P_{c,\lambda}^{(\le P)}-P_{c,\lambda}
 +q_X\log b_\delta
 +\frac{q_X}{2}\log\left(1+\frac{R_{P(X)}(X,c)}{q_X\delta^2}\right)
 \nonumber\\
 &+T_X+\frac{C_{\rm out}+18}{L_X}
       +\frac{16}{3\pi\kappa_E}P^{3/8}X^{-4}.
\end{align}
Here \(q_X\log b_\delta\le\delta/4\). All comparison costs can
vanish with a suitable regulator; the exact arithmetic residual in
\eqref{eq:exact-relative-residual} is still unpaid. In particular
\eqref{eq:adaptive-provider-error} estimates \(V-S_\lambda\),
not \(V-S_P^{\rm ref}\). Neither the fixed-\(P\) mean law nor
the fixed-\(P\) unboundedness theorem permits substituting this
growing cutoff into its quantifiers.

\subsection{The specified signed tangent still fails for fixed \(P\)}
Retaining a signed cross term is not by itself enough. Fix finite
\(P\), take \(c=1\), and set
\[
 w=V/E_P,\quad a=S_P^{\rm ref}/E_P,\quad h=|a|^2+\delta^2,
 \qquad G_\delta=\tfrac12\log h-\tfrac12+\frac{|w|^2}{2h}.
\]
The inequality \(\log t\le t-1\) gives \(\log|w|\le G_\delta\),
with the exact cross term retained in
\(|w|^2=|a|^2+2\Re(\bar a(w-a))+|w-a|^2\).
Let \(H=4+\delta^2\) and \(Z=\int_{W_X}|w|^2\,d\omega_X\).
Since \(|a|\le2\), \(h\le H\) and \(h/H\ge|a|^2/4\).
Integrating, and then using \(t-1\ge\log t\) at \(t=Z/(q_XH)\),
gives the regulator-independent lower bound
\begin{align}\label{eq:tangent-obstruction}
 \int_{W_X}G_\delta\,d\omega_X
 &\ge\frac{q_X}{2}\log H-\frac{q_X}{2}+\frac Z{2H}
         +\int_{W_X}\log|a|\,d\omega_X-q_X\log2\nonumber\\
 &\ge\frac{q_X}{2}\log(Z/q_X)
         +\int_{W_X}\log|a|\,d\omega_X-q_X\log2.
\end{align}
The reference logarithm is \(O_P(1/L_X)\) by
Lemma~\ref{lem:joint-phase}, including its zeros. For fixed \(P\),
\(Z\ge M(X,1)/B_P^2\). Theorem~\ref{thm:fixed-residual-unbounded}
therefore makes the integrated tangent unbounded at arbitrarily high
centers, independently of \(\delta\). Restoring the bounded reference
prime difference and the nonnegative tail allowances preserves this
failure. The result concerns this fixed-reference tangent, with a
regulator constant on each window. It does not cover a growing
reference, spatially varying tangent heights, or all nonlinear signed
estimates, and gives no lower bound for \(J\).
